\documentclass{article}
\usepackage{geometry}
\geometry{a4paper, margin=1in}
\usepackage{hyperref}
\hypersetup{colorlinks=true, urlcolor=blue, citecolor=blue}

\title{Proposal A: Self-Supervised Feature Pretraining for\\Label-Efficient Cell Tracking in Embryogenesis}
\author{Candidate Name}
\date{Suggested Supervisor: Prof.\ Dr.\ Martin Weigert}

\begin{document}
\maketitle

\section*{Motivation}

Trackastra~\cite{trackastra} is a transformer-based cell tracking method that links pre-segmented cells across time frames by predicting pairwise association scores from image crops. It won the 7th Cell Tracking Challenge (ISBI 2024) and ships with pretrained models for 2D cell culture data. However, applying Trackastra to \textbf{embryogenesis data}---3D light-sheet time-lapses of developing organisms such as \textit{Tribolium castaneum}---faces a domain gap: the pretrained encoder has never seen the dense, noisy, tissue-scale fluorescence images characteristic of embryo recordings. Bridging this gap via supervised fine-tuning would require expensive manual track annotations in 3D, which is precisely the bottleneck Trackastra was designed to alleviate.

A recent variant of Trackastra replaces raw image crops with features from SAM2~\cite{sam2}, a foundation model pretrained on natural images and video. This shows the encoder is modular and that better features improve tracking. But SAM2 features---like DINOv2---are learned from natural photographs and carry a significant domain gap to fluorescence microscopy. The Weigert group's own TAP~\cite{tap} (Time Arrow Prediction) learns dense representations from \textbf{unlabeled} live-cell video by predicting temporal order---an SSL objective specifically designed for microscopy time-lapse data. TAP features have been shown to capture biologically meaningful events (cell division, migration) without any labels.

\textbf{These two tools have never been connected.} This project bridges TAP and Trackastra by using self-supervised temporal features as the input representation for cell tracking, evaluated on embryogenesis data.

\section*{Approach}

\begin{enumerate}
	\item \textbf{TAP pretraining on unlabeled embryo data.} Train TAP's U-Net encoder on raw (unlabeled) time-lapse frames from the SLICE-1 dataset~\cite{slice1} (50 \textit{Tribolium} embryos, 5 transgenic lines, 5.7\,M images) or the Cell Tracking Challenge \textit{Tribolium} sequence (Fluo-N3DL-TRIC). No annotations are required---only the raw video.
	\item \textbf{Feature extraction and integration.} Extract dense per-frame feature maps from the TAP-pretrained encoder. For each segmented cell, pool features from its mask region to obtain a cell-level embedding. Feed these embeddings into Trackastra's association transformer in place of the default raw-crop encoder.
	\item \textbf{Systematic comparison of feature backbones.} Evaluate tracking quality with:
	      \begin{itemize}
		      \item Trackastra default (raw crops, trained from scratch)
		      \item Trackastra + SAM2 features (existing variant)
		      \item Trackastra + frozen DINOv2 features (natural-image SSL)
		      \item Trackastra + TAP features (microscopy-specific temporal SSL)
		      \item Trackastra + TAP features with LoRA fine-tuning on limited annotations
	      \end{itemize}
	\item \textbf{Evaluation on CTC benchmarks and \textit{Tribolium} data.} Use the established Cell Tracking Challenge metrics (TRA, DET, AOGM) for quantitative comparison. Qualitatively inspect tracked lineage trees on \textit{Tribolium} embryo data.
\end{enumerate}

\section*{Why This Works as a One-Semester Project}

\begin{itemize}
	\item Both TAP and Trackastra are \textbf{existing PyTorch codebases} from the Weigert group. The core engineering task is integration, not invention.
	\item The CTC provides \textbf{established benchmarks} with ground truth, so evaluation is well-defined.
	\item SLICE-1 provides \textbf{abundant unlabeled embryo data} for SSL pretraining without annotation effort.
	\item The comparison matrix (5 feature variants) produces \textbf{clean, publishable results} regardless of which method wins.
	\item The project is \textbf{incrementally buildable}: TAP pretraining works first, Trackastra integration second, evaluation third.
\end{itemize}

\section*{Expected Contribution}

A quantitative answer to the question: \emph{Can self-supervised temporal features, learned from unlabeled microscopy video, improve cell tracking---and can they outperform general-purpose foundation model features (SAM2, DINOv2) that were learned from natural images?} If yes, this establishes SSL pretraining as a practical strategy for label-efficient tracking in embryogenesis. If no, the comparison itself is informative and identifies where the feature gap lies.

\begin{thebibliography}{9}
	\bibitem{trackastra} B.~Gallusser and M.~Weigert. ``Trackastra: Transformer-based cell tracking for live-cell microscopy.'' \textit{ECCV}, 2024.
	\bibitem{tap} B.~Gallusser, M.~Stieber, and M.~Weigert. ``Self-supervised dense representation learning for live-cell microscopy with time arrow prediction.'' \textit{MICCAI}, 2023.
	\bibitem{sam2} N.~Ravi et al. ``SAM 2: Segment Anything in Images and Videos.'' \textit{arXiv:2408.00714}, 2024.
	\bibitem{slice1} F.~Strobl et al. ``SLICE-1: A comprehensive light-sheet imaging dataset of \textit{Tribolium castaneum} embryogenesis.'' \textit{Scientific Data}, Dec 2025.
\end{thebibliography}

\end{document}
